\section{Introduction}
Width-based search algorithms have in the last few years become among the state-of-the-art approaches to  automated planning, e.g.\ the original Iterated Width (IW) algorithm~\cite{iworginal}. As in propositional STRIPS planning, states are represented by a set of propositional literals, also called boolean \emph{features}. The state space is searched with breadth-first search (BFS), but the state space explosion problem is handled by pruning states based on their \emph{novelty}. First, a parameter $k$ is chosen, called the \emph{width parameter} of the search. Searching with width parameter $k$ essentially means that we only consider $k$ literals/features at a time. A state $s$ generated during the search is called \emph{novel} if there exists a set of $k$ literals/features not made true in any earlier generated state. Unless a state is novel, it is immediately pruned from the search. Clearly, we then reduce the size of the searched state space to be exponential in $k$. It has been shown that many classical planning problems, e.g.\ problems from the International Planning Competition (IPC) domains, can be solved efficiently using width-based search with very low values of $k$.  

The essential benefit of using width-based algorithms is the ability to perform semi-structured (based on feature structures) exploration of the state space, and reach deep states that may be important for achieving the planning goals. In classical planning, width-based search has been integrated with heuristic search methods, leading to Best-First Width Search~\cite{BFSW} that performed well at the 2017 IPC. Width-based search has also been adapted to reward-driven problems where the algorithm uses a simulator for interacting with the environment \cite{ijcai2017-600}. This has enabled the use of width-based search in reinforcement learning environments such as the Atari 2600 video game suite, through the Arcade Learning Environment (ALE)~\cite{DBLP:journals/corr/abs-1207-4708}. There have been several implementations of width-based search directly using the RAM states of the Atari computer as features \cite{ram_model,ram_IW_shleyfman,dominatedSequnce}.

Motivated by the fact that humans do not have access to RAM states when playing video games, methods for using these algorithms from screen pixels have been developed. \citet{rolloutiw-orginal} propose a modified version of IW, called RolloutIW, that uses pixel-based features and achieves results comparable to learning methods in almost real-time on ALE. The combination of reinforcement learning and RolloutIW in the pixel domain is explored by \citet{RLIW}, who propose to train a policy to guide the action selection in rollouts. 

% Regardless of different efforts, 
Width-based algorithms are highly exploratory and, not surprisingly, significantly dependent on the quality of the features defined. 
% The original RolloutIW paper~\cite{rolloutiw-orginal} uses 
\citet{rolloutiw-orginal} use
a set of low-level, color-based features called B-PROST~\cite{liang},
designed by humans to achieve good performance in ALE. Integrating feature learning
into the algorithms themselves rather than using hand-crafted features is an interesting challenge---and a challenge that will bring the algorithms more on a par with humans learning to play those video games. The features used by \citet{RLIW} were the output features of the last layer of the policy network, improving the performance of IW. With the attempt to learn efficient features for width-based search, and inspired by the results of \citet{RLIW}, this paper investigates the possibility of a more structured approach to generate features for width-based search using deep generative models. 

More specifically, in this paper we use variational autoencoders (VAE)---latent variable models that have been widely used for representation learning \cite{kingma2019introduction} as they allow for approximate inference of the latent variables underlying the data---to learn propositional symbols directly from pixels and without any supervision.
We then investigate whether the learned symbolic representation can be successfully used for planning in the Atari 2600 domain. 

We compare the planning performance of RolloutIW using our learned representations with RolloutIW using the handcrafted B-PROST features and report human scores as a baseline.
Our results show that our learned representations lead to better performance than B-PROST on RolloutIW (and often outperform human players), although the B-PROST features also contain temporal information (how the screen changes), whereas we only extract static features from individual frames. 
Apart from improving scores in the Atari 2600 domain, we also significantly reduce the number of features (by a factor of more than $10^3$). 
We also investigate in detail (with ablation studies) which factors seem to lead to good performance on games in the Atari domain.

The main contributions of our paper are hence:
\begin{itemize}
    \item We propose a novel combination of variational autoencoders and width-based planning algorithms.
    \item We train VAEs to learn propositional symbols for planning, from pixels and without any supervision.
    \item We run a large-scale evaluation on Atari 2600 where we compare the performance of RolloutIW with our learned representations to RolloutIW with B-PROST features.
    \item We investigate with ablation studies the effect of various hyperparameters in both learning and planning.
    \item We show that our learned representations lead to higher scores than using B-PROST, even though (i) our features don't have any temporal information and (ii) our models are trained from data collected by RolloutIW with B-PROST, thus limiting the richness of the dataset.
\end{itemize}
Below, we provide the required background on RolloutIW, B-PROST, and VAEs, present the original approach of this paper, discuss our experimental results, and conclude.
